\begin{abstract}
Pattern matching is a foundation of functional and dependently typed languages, yet it often forces programmers to handle clauses that can never be reached.
Programmers must then either write unsafe partial functions or, in safe settings, discharge tedious proof obligations for these trivial cases.
The problem is particularly critical when programs operate on subsets of an inductive datatype.
Standard encodings pair data with propositions or refactor datatypes into indexed families, but both approaches complicate composition and proof automation.

The problem has two causes:
(1)~the type system cannot express the intended subset of the data, so it generates trivial obligations, and
(2)~it cannot reason computationally about constructors, so ultimately no procedure can be written to discharge them.
A subtyping discipline addresses the first; the second requires lifting constructors to the term level, which is more expressive and can itself encode the subtyping discipline.
In fact, the latter approach can also be used to encode other procedures that eliminate additional trivial proof obligations.

We present a subtyping discipline for \emph{constructor subsets}, where first-class signatures track which constructors of a family are available, so that pattern matching rules out impossible branches without spurious obligations.
We elaborate it into the type theory of Dagand and McBride, in which constructors are themselves terms, and we prove that our elaboration is sound.

Our results show that constructor subsets are expressive enough to capture common programming idioms, and that extending a dependently typed language with them is simple.
Through the elaboration, constructor subsets can be expressed in a purely computational setting, which enables reusing constructors in different type constructors, deciding about constructors, and even set-based operations on signatures.
We implement the discipline in Lambdapi and mechanize the elaborated calculus in Rocq.
\end{abstract}
